% Poème grothendieckien.
% Note courte, dans la manière des Éléments (site, champ, gerbe, descente, point de vue relatif).
% Le vocabulaire reste abstrait à dessein : la nature des deux problèmes se devine en lisant.
\documentclass[10pt]{article}
\usepackage[a4paper,left=42mm,right=42mm,top=24mm,bottom=26mm]{geometry}
\usepackage{fontspec}
\usepackage{amsmath,amssymb,amsthm}
\usepackage[french]{babel}
\usepackage[hidelinks]{hyperref}
\usepackage{enumitem}
\setlist{nosep,leftmargin=1.4em}
\setlength{\parskip}{0.25em}
\setlength{\emergencystretch}{2em}
\frenchsetup{SmallCapsFigTabCaptions=false}

\newtheoremstyle{bbk}{0.45em}{0.45em}{\itshape}{}{\bfseries}{.}{0.4em}{}
\theoremstyle{bbk}
\newtheorem{prop}{Proposition}[section]
\newtheorem{coro}[prop]{Corollaire}
\newtheorem{conj}[prop]{Conjecture}
\newtheoremstyle{bbkdef}{0.45em}{0.45em}{}{}{\bfseries}{.}{0.4em}{}
\theoremstyle{bbkdef}
\newtheorem{defi}[prop]{Définition}
\newtheorem{rema}[prop]{Remarque}
\renewcommand{\qedsymbol}{\(\square\)}

\newcommand{\Lec}{\mathcal L}
\newcommand{\U}{\mathfrak U}
\newcommand{\Cl}[1]{\overline{#1}}
% Un mot caviardé : un rectangle noir de sa largeur ; le mot lui-même n'entre pas dans le PDF.
\newcommand{\cache}[1]{\setbox0\hbox{#1}\rule[-0.32ex]{\wd0}{1.02em}}
\makeatletter
\renewcommand{\@seccntformat}[1]{\S\,\csname the#1\endcsname.\quad}
\makeatother

\title{\vspace{-1.2em}\large\bfseries Poème grothendieckien\\[0.2em]
\normalsize\mdseries Deux problèmes de recollement}
\author{\small Michel Hua\thanks{Rédigé avec un assistant, Claude Opus~5.5 (Anthropic), puis relu par l'auteur.}}
\date{\small Novembre 2026}

\begin{document}
\maketitle
\vspace{-1.6em}

\noindent\emph{Introduction.} On se donne des fragments, chacun pris dans sa propre langue. On
les traite comme les objets locaux d'un champ ; ce qu'ils esquissent ensemble est une section
globale, et ce qu'ils résolvent se lit sur une base. Les obstructions sont des classes de
cohomologie, et c'est elles qu'il s'agit de lire.

\section{Le champ des lectures}

\begin{defi}
Soit $P$ un ensemble fini de \emph{fragments}, et $S$ un ensemble de symboles. Une partie
$U \subset P$ est \emph{autonome} si elle porte le contexte qu'il faut pour être comprise seule.
Les parties autonomes et leurs inclusions forment une catégorie $C$, que l'on munit de la
topologie engendrée par les familles $(U_i \subset U)$ de réunion $U$. Un fragment isolé n'est en
général pas autonome : c'est ce qui rend cette topologie non triviale.
\end{defi}

\begin{defi}
Le \emph{champ des lectures} $\Lec$ associe à $U$ le groupoïde $\Lec(U)$ dont les objets sont
les lectures de $U$ — un sens donné à chaque fragment, et une application qui envoie chaque
symbole employé dans $U$ sur l'objet qu'il désigne — et dont les flèches $\ell \to \ell'$ sont les
\emph{renommages} des symboles employés qui transforment $\ell$ en $\ell'$. La restriction à
$V \subset U$ oublie les fragments hors de $V$.
\end{defi}

\begin{prop}
$\Lec$ est un préchamp : pour $\ell, \ell' \in \Lec(U)$, le préfaisceau
$V \mapsto \mathrm{Isom}(\ell|_V, \ell'|_V)$ est un faisceau. Une donnée de descente
$(\ell_i, g_{ij})$ relative à un recouvrement $\U$ — lectures $\ell_i \in \Lec(U_i)$ et
renommages $g_{ij} \colon \ell_j|_{U_{ij}} \to \ell_i|_{U_{ij}}$ tels que
$g_{ij}g_{jk} = g_{ik}$ — est effective, c'est-à-dire provient d'une lecture de $U$ dans une
notation unique, si et seulement si sa classe est nulle dans $\check H^1(\U, G)$, où $G$ est le
groupe des permutations de $S$.
\end{prop}

\begin{proof}
Des renommages compatibles sur un recouvrement se recollent, un renommage n'étant qu'une
application définie symbole par symbole. Si $g_{ij} = h_i^{-1}h_j$, les lectures $h_i\ell_i$
coïncident sur les intersections et se recollent ; la réciproque est la définition de la
trivialité de la classe.
\end{proof}

Le champ associé $\tilde\Lec$ a donc plus d'objets que $\Lec$ : ce sont les \emph{textes},
données de descente que l'on accepte de ne pas réduire à une notation unique. Un objet de
$\tilde\Lec(P)$ est un texte de $P$ ; un objet de $\Lec(P)$ en est une \emph{réalisation}.

\begin{coro}
Si le nerf de $\U$ est un arbre, tout texte relatif à $\U$ admet une réalisation.
\end{coro}

\begin{proof}
$\check H^1$ d'un arbre est trivial à coefficients dans un groupe quelconque : on choisit $h$ en
une racine et on le propage le long des arêtes par $h_j = h_ig_{ij}$.
\end{proof}

\begin{prop}
Supposons la lecture \emph{localement bien posée} : tout fragment appartient à une partie
autonome qui admet une lecture, et deux lectures d'une même partie sont localement isomorphes.
Alors $\tilde\Lec$ est une gerbe ; son lien est le faisceau $\mathcal A$ des automorphismes d'une
lecture, les \emph{symétries de notation}. Un texte de $P$ existe si et seulement si cette gerbe
est neutre ; lorsque $\mathcal A$ est abélien, l'obstruction est une classe de
$H^2(C, \mathcal A)$, et deux textes diffèrent par un torseur sous $\mathcal A$, classé par
$H^1(C, \mathcal A)$ (Giraud).
\end{prop}

\begin{proof}
Les deux hypothèses sont, mot pour mot, les axiomes d'une gerbe ; le reste est la théorie de
Giraud des gerbes à lien abélien.
\end{proof}

\begin{rema}
On voit ainsi deux étages d'obstruction. Le premier ($H^2$) dit qu'il n'y a pas de texte du tout :
les symétries locales ne se recollent pas. Le second ($H^1$) dit qu'il y a un texte, mais pas de
réalisation dans une notation unique. Une classe non nulle n'est pas une faute : elle naît d'un
cycle de fragments où un même nom désigne deux objets, et on la tue en agrandissant $S$, non en
corrigeant une lecture.
\end{rema}

\begin{defi}
Un fragment peut admettre plusieurs lectures. Soit $\Lambda(U) \subset \mathrm{Ob}\,\Lec(U)$
l'ensemble des lectures \emph{admissibles} de $U$, et $E(\lambda)$ l'ensemble des énoncés que
$\lambda$ affirme. Un énoncé est \emph{reconstitué} sur $U$ s'il appartient à
\[E(U) = \textstyle\bigcap_{\lambda \in \Lambda(U)} E(\lambda),\]
et \emph{correct} s'il est vrai tel qu'il est écrit, éventuellement après l'ajout d'une hypothèse
que le fragment tait et que l'on déclare.
\end{defi}

Restreindre $\Lambda$ fait croître $E$, et un énoncé reconstitué et correct le reste. Une
vérification \cache{formelle} contrôle donc les mathématiques d'une lecture, jamais sa fidélité au
fragment : les deux se mesurent séparément.

\section{Le point de vue relatif}

Tout ce qui précède se place au-dessus d'une base : un ensemble ordonné $T$, muni de sa topologie
d'Alexandrov, que l'on appellera le \cache{temps}. Les fragments y sont placés par un morphisme
$P \to T$ ; on note $t_P$ la borne de son image. Au-dessus de $T$ vit un faisceau croissant
$t \mapsto C_{<t}$, les énoncés que la communauté \cache{mathématique} tenait pour acquis avant $t$, et un
ensemble $Q$ de \cache{questions}, chacune avec sa date $t_q$. Un ordre des fragments compatible aux
indices de position existe si et seulement si ces indices sont sans cycle (Szpilrajn).

\begin{prop}
Soit $R \subset Q \times E$ la relation « l'énoncé $e$ touche la question $q$ ». Les applications
$A \mapsto A'$ et $B \mapsto B'$ qu'elle définit entre parties de $Q$ et de $E$ forment une
correspondance de Galois, et les couples fermés un treillis complet.
\end{prop}

\begin{proof}
C'est la polarité de Birkhoff d'une relation binaire ; les couples fermés sont les « concepts »
de Ganter et Wille.
\end{proof}

\begin{defi}
Soit $\vdash$ un système déductif où chaque question fermée $q$ s'énonce par une formule
$\varphi_q$, et $\Cl X$ la clôture déductive de $X$. Une \emph{réponse} à $q$ est une partie finie
$B \subset E(P)$ telle que $\varphi_q$ ou $\neg\varphi_q$ appartienne à $\Cl B$ ; elle est
\emph{attestée} si chaque énoncé de $B$ porte son fragment, \emph{vérifiée} si la déduction a été
contrôlée, \emph{anticipée} si tous ses fragments précèdent $t_q$.
\end{defi}

\begin{prop}
Si $\varphi_q \in \Cl{E(P)}$, il existe une réponse minimale pour l'inclusion, et l'ensemble des
questions résolues croît avec $P$.
\end{prop}

\begin{proof}
Une démonstration n'emploie qu'un nombre fini d'hypothèses.
\end{proof}

\begin{rema}
Une réponse anticipée est une section au-dessus de l'ouvert $\{t < t_q\}$, qui ne contient pas la
date de la question : la question avait une réponse avant d'être posée. La donner après coup
n'établit aucune priorité ; cela fournit, avec les fragments, la matière d'une \cache{histoire}. Une
question peut aussi n'avoir de réponse que dans le texte recollé, et dans aucune lecture locale.
\end{rema}

Soit enfin $D(P) \subset E(P)$ l'ensemble des énoncés dont les fragments donnent eux-mêmes une
démonstration.

\begin{conj}
Tout énoncé correct de $E(P)$ appartient à
\[\Cl{D(P) \cup C_{<t_P}} .\]
L'espace des énoncés démontrables que visent les fragments se reconstitue à partir d'eux et de
ce que la communauté savait alors.
\end{conj}

\begin{rema}
Un contre-exemple serait un énoncé correct que les fragments affirment sans le démontrer et que
rien, en leur temps, ne démontrait : une anticipation. C'est en cherchant les contre-exemples
qu'on trouve ce que $P$ a de neuf ; chaque réponse anticipée qui tombe hors de
$\Cl{D(P) \cup C_{<t_P}}$ en est un.
\end{rema}

\begin{rema}[Réflexion]
On aurait pu attaquer au marteau : fixer d'avance une notation et y plier chaque fragment. Le
yoga est ici inverse. On laisse chaque fragment dans sa langue, on note seulement comment les
langues se traduisent sur les recouvrements, et l'on attend que la structure parle : la
cohomologie dit si un texte existe et s'il admet une réalisation. Quand elle ne s'annule pas, elle
ne dit pas que la lecture est mauvaise ; elle dit où regarder.
\end{rema}

\vspace{0.3em}
\noindent{\small\textbf{Références.}
G.~Birkhoff, \emph{Lattice Theory}, Amer. Math. Soc. Colloq. Publ.~25, 1940. ---
B.~Ganter, R.~Wille, \emph{Formal Concept Analysis. Mathematical Foundations}, Springer, 1999. ---
J.~Giraud, \emph{Cohomologie non abélienne}, Grundlehren~179, Springer, 1971. ---
R.~Godement, \emph{Topologie algébrique et théorie des faisceaux}, Hermann, 1958. ---
E.~Szpilrajn, « Sur l'extension de l'ordre partiel », \emph{Fund. Math.} 16 (1930), 386--389.}

\end{document}
